\section{Sujet n\textsuperscript{o}\,11}
\noindent\begin{minipage}{0.5\linewidth}\footnotesize Office du Baccalauréat du Cameroun\end{minipage}%
\hfill\begin{minipage}{0.45\linewidth}\footnotesize\raggedleft Examen : \textbf{Baccalauréat} · Série : \textbf{C/E}\\
Épreuve : \textbf{Mathématiques} · Durée : \textbf{4\,h} · Coef. : \textbf{7}\end{minipage}
\par\smallskip{\color{onbo}\rule{\linewidth}{1pt}}

\subsection*{Partie A — Évaluation des ressources \hfill (15 points)}

\noindent\textbf{Exercice 1.} \hfill\textit{(3 points)}\\
Une enquête menée par un opérateur de téléphonie montre que $60\,\%$ de ses abonnés ont
choisi le forfait « Internet ». Parmi les abonnés au forfait Internet, $70\,\%$ se déclarent
satisfaits ; parmi les autres abonnés, $40\,\%$ se déclarent satisfaits. On interroge au
hasard un abonné. On note $I$ l'événement « l'abonné a le forfait Internet » et $S$
l'événement « l'abonné est satisfait ».
\begin{enumerate}[label=\arabic*)]
\item Construire un arbre pondéré décrivant la situation. \hfill\textit{(0,75 pt)}
\item Calculer la probabilité que l'abonné ait le forfait Internet et soit satisfait. \hfill\textit{(0,75 pt)}
\item Montrer que $P(S)=0{,}58$. \hfill\textit{(0,75 pt)}
\item L'abonné interrogé est satisfait. Quelle est la probabilité qu'il ait le forfait Internet ? \hfill\textit{(0,75 pt)}
\end{enumerate}

\smallskip
\noindent\textbf{Exercice 2.} \hfill\textit{(3 points)}\\
On considère la suite $(u_n)$ définie par $u_0=5$ et, pour tout $n\in\N$,
$u_{n+1}=\dfrac12 u_n+3$.
\begin{enumerate}[label=\arabic*)]
\item On pose $v_n=u_n-6$. Montrer que $(v_n)$ est géométrique ; préciser sa raison et $v_0$. \hfill\textit{(1,5 pt)}
\item Exprimer $u_n$ en fonction de $n$, puis déterminer $\displaystyle\lim_{n\to+\infty}u_n$. \hfill\textit{(1,5 pt)}
\end{enumerate}

\smallskip
\noindent\textbf{Exercice 3.} \hfill\textit{(4 points)}\\
Soit $f$ la fonction définie sur $\R$ par $f(x)=(2-x)\mathrm{e}^{x}$, de courbe $(\mathcal{C})$.
\begin{enumerate}[label=\arabic*)]
\item Déterminer les limites de $f$ en $-\infty$ et $+\infty$ ; interpréter graphiquement
celle en $-\infty$. \hfill\textit{(1 pt)}
\item Montrer que $f'(x)=(1-x)\mathrm{e}^{x}$ et dresser le tableau de variations de $f$. \hfill\textit{(1,5 pt)}
\item À l'aide d'une intégration par parties, calculer $\displaystyle\int_0^{1}f(x)\,\dd x$. \hfill\textit{(1,5 pt)}
\end{enumerate}

\smallskip
\noindent\textbf{Exercice 4.} \hfill\textit{(5 points)}\\
L'espace est muni d'un repère orthonormé $(O\,;\,\vec\imath,\vec\jmath,\vec k)$. On considère
les points $A(2\,;\,0\,;\,1)$, $B(1\,;\,2\,;\,3)$ et $C(0\,;\,1\,;\,1)$.
\begin{enumerate}[label=\arabic*)]
\item Démontrer que les points $A$, $B$, $C$ ne sont pas alignés. \hfill\textit{(1 pt)}
\item Montrer que $\vec n(2\,;\,4\,;\,-3)$ est un vecteur normal au plan $(ABC)$,
puis donner une équation cartésienne de $(ABC)$. \hfill\textit{(1,5 pt)}
\item Déterminer une représentation paramétrique de la droite $(\Delta)$ passant par
$O$ et perpendiculaire au plan $(ABC)$. \hfill\textit{(1 pt)}
\item Calculer la distance du point $O$ au plan $(ABC)$. \hfill\textit{(1,5 pt)}
\end{enumerate}

\subsection*{Partie B — Évaluation des compétences \hfill (5 points)}

\noindent\textit{La société « CamTel+ » étudie le déploiement de ses abonnés. Elle te sollicite,
en tant qu'élève de Terminale~C, pour trois analyses. Le nombre d'abonnés à un nouveau
forfait est de $40$ milliers le premier mois, puis chaque mois la société conserve
$80\,\%$ de ses abonnés et en gagne $15$ milliers de nouveaux. Par ailleurs, le responsable
veut générer des codes promotionnels formés de $3$ lettres distinctes (parmi $26$) suivies
de $2$ chiffres distincts (parmi $10$). Enfin, un test de couverture réseau est positif pour
$5\,\%$ des zones réellement mal couvertes ; on contrôle $6$ zones mal couvertes au hasard.}

\begin{enumerate}[label=\textbf{Tâche \arabic*.},leftmargin=2.4em]
\item En modélisant le nombre d'abonnés (en milliers) du mois $n$ par une suite $(a_n)$
avec $a_1=40$ et $a_{n+1}=0{,}8\,a_n+15$, déterminer la limite de $(a_n)$ et l'interpréter. \hfill\textit{(1,5 pt)}
\item Déterminer le nombre de codes promotionnels distincts que la société peut générer. \hfill\textit{(1,5 pt)}
\item Déterminer la probabilité qu'\emph{au moins une} des $6$ zones contrôlées soit
détectée comme mal couverte par le test. \hfill\textit{(1,5 pt)}
\end{enumerate}
\par\smallskip{\footnotesize\itshape Présentation générale : 0,5 point.}

% ════════════════════════════════════════════════════
\corrige{du sujet 11}

\noindent\textbf{Partie A — Exercice 1.}
1) Arbre : depuis la racine, $P(I)=0{,}6$ et $P(\bar I)=0{,}4$ ; sous $I$ : $P_I(S)=0{,}7$,
$P_I(\bar S)=0{,}3$ ; sous $\bar I$ : $P_{\bar I}(S)=0{,}4$, $P_{\bar I}(\bar S)=0{,}6$.
2) $P(I\cap S)=P(I)\times P_I(S)=0{,}6\times0{,}7=0{,}42$.
3) Formule des probabilités totales : $P(S)=P(I\cap S)+P(\bar I\cap S)=0{,}42+0{,}4\times0{,}4=0{,}42+0{,}16=0{,}58$.
4) $P_S(I)=\dfrac{P(I\cap S)}{P(S)}=\dfrac{0{,}42}{0{,}58}=\dfrac{42}{58}=\dfrac{21}{29}\approx0{,}724$.

\noindent\textbf{Exercice 2.}
1) $v_{n+1}=u_{n+1}-6=\tfrac12u_n+3-6=\tfrac12u_n-3=\tfrac12(u_n-6)=\tfrac12v_n$. Donc $(v_n)$
est géométrique de raison $q=\tfrac12$ et de premier terme $v_0=u_0-6=-1$.
2) $v_n=-\big(\tfrac12\big)^n$, d'où $u_n=v_n+6=6-\big(\tfrac12\big)^n$. Comme
$\big(\tfrac12\big)^n\to0$, $\displaystyle\lim_{n\to+\infty}u_n=6$.

\noindent\textbf{Exercice 3.}
1) En $+\infty$ : $2-x\to-\infty$ et $\mathrm e^x\to+\infty$, donc $f(x)\to-\infty$.
En $-\infty$ : posons-le par croissance comparée, $x\mathrm e^x\to0$ et $\mathrm e^x\to0$,
donc $f(x)=2\mathrm e^x-x\mathrm e^x\to0$ : la droite $y=0$ est asymptote horizontale en $-\infty$.
2) $f'(x)=(-1)\mathrm e^x+(2-x)\mathrm e^x=(1-x)\mathrm e^x$. Le signe de $f'$ est celui de
$1-x$ : $f$ est croissante sur $]-\infty,1]$, décroissante sur $[1,+\infty[$, avec un
maximum $f(1)=(2-1)\mathrm e^1=\mathrm e$.
3) IPP avec $u=2-x$, $u'=-1$, $v'=\mathrm e^x$, $v=\mathrm e^x$ :
$\int_0^1(2-x)\mathrm e^x\dd x=\big[(2-x)\mathrm e^x\big]_0^1+\int_0^1\mathrm e^x\dd x
=\big[(2-x)\mathrm e^x\big]_0^1+\big[\mathrm e^x\big]_0^1$.
Or $\big[(2-x)\mathrm e^x\big]_0^1=1\cdot\mathrm e-2\cdot1=\mathrm e-2$ et $\big[\mathrm e^x\big]_0^1=\mathrm e-1$,
donc $\int_0^1f(x)\dd x=(\mathrm e-2)+(\mathrm e-1)=2\mathrm e-3$.

\noindent\textbf{Exercice 4.}
1) $\vec{AB}(-1\,;\,2\,;\,2)$ et $\vec{AC}(-2\,;\,1\,;\,0)$. Ces vecteurs ne sont pas
colinéaires (par exemple $-1/-2\ne2/1$), donc $A$, $B$, $C$ ne sont pas alignés.
2) $\vec n\cdot\vec{AB}=2(-1)+4(2)+(-3)(2)=-2+8-6=0$ et
$\vec n\cdot\vec{AC}=2(-2)+4(1)+(-3)(0)=-4+4+0=0$. Le vecteur $\vec n$ est orthogonal à
deux vecteurs directeurs non colinéaires du plan : c'est un vecteur normal à $(ABC)$.
Équation : $2x+4y-3z+d=0$ ; avec $A(2\,;\,0\,;\,1)$, $4+0-3+d=0$ donne $d=-1$. D'où
$(ABC):\ 2x+4y-3z-1=0$.
3) $(\Delta)$ passe par $O$ et a pour vecteur directeur $\vec n(2\,;\,4\,;\,-3)$ :
$\begin{cases}x=2t\\ y=4t\\ z=-3t\end{cases}\ (t\in\R)$.
4) $d(O,(ABC))=\dfrac{|2\cdot0+4\cdot0-3\cdot0-1|}{\sqrt{2^2+4^2+(-3)^2}}=\dfrac{1}{\sqrt{4+16+9}}=\dfrac{1}{\sqrt{29}}=\dfrac{\sqrt{29}}{29}\approx0{,}186$.

\smallskip
\noindent\textbf{Partie B — Situation.}
\emph{Tâche 1.} La suite vérifie $a_{n+1}=0{,}8\,a_n+15$. On cherche le point fixe $\ell$ :
$\ell=0{,}8\ell+15\Rightarrow0{,}2\ell=15\Rightarrow\ell=75$. En posant $w_n=a_n-75$, on a
$w_{n+1}=0{,}8\,w_n$ (géométrique, raison $0{,}8\in\,]0,1[$), donc $w_n\to0$ et
$a_n\to75$. À long terme, le nombre d'abonnés se stabilise autour de
\textbf{75 milliers (75\,000 abonnés)}.

\emph{Tâche 2.} Choix ordonné de $3$ lettres distinctes : arrangements $A_{26}^3=26\times25\times24=15\,600$.
Choix ordonné de $2$ chiffres distincts : $A_{10}^2=10\times9=90$. Comme la position
lettres/chiffres est fixée, le nombre de codes est $15\,600\times90=\mathbf{1\,404\,000}$.

\emph{Tâche 3.} Chaque zone mal couverte est détectée par le test avec probabilité $0{,}05$,
les $6$ contrôles étant indépendants. Donc
$P(\text{au moins une détectée})=1-(0{,}95)^6\approx1-0{,}7351=\mathbf{0{,}2649}$, soit environ \SI{26,5}{\percent}.
